\section{Subsequences}

\import{set/sequence-filter.tex}

A subsequence is obtained by precomposing a sequence with a strictly
increasing sequence of natural indices.  Keeping the selector explicit makes
the order content reusable independently of any topology on the codomain.

\begin{definition}\label{natural_subsequence_selector}
    $h$ is a natural subsequence selector iff
        $h\in\funs{\naturals}{\naturals}$ and for all
        $m,n\in\naturals$ we have if $m\in n$, then $h(m)\in h(n)$.
\end{definition}

\begin{proposition}\label{natural_subsequence_selector_type}
    Suppose $h$ is a natural subsequence selector.
    Then $h\in\funs{\naturals}{\naturals}$.
\end{proposition}
\begin{proof}
    Follows by \obv{natural_subsequence_selector}.
\end{proof}

\begin{proposition}\label{natural_subsequence_selector_strict}
    Suppose $h$ is a natural subsequence selector.
    Let $m,n\in\naturals$.
    Suppose $m\in n$.
    Then $h(m)\in h(n)$.
\end{proposition}
\begin{proof}
    Follows by \obv{natural_subsequence_selector}.
\end{proof}

\begin{proposition}\label{natural_subsequence_selector_successor}
    Suppose $h$ is a natural subsequence selector.
    Let $n$ be a natural number.
    Then $h(n)\in h(\suc{n})$.
\end{proposition}
\begin{proof}
    $\suc{n}$ is a natural number by \obv{omega_successor}.
    $n\in\suc{n}$ by \obv{suc_intro_self}.
    Follows by \obv{natural_subsequence_selector_strict}.
\end{proof}

\begin{theorem}\label{natural_subsequence_selector_from_successors}
    Let $h\in\funs{\naturals}{\naturals}$.
    Suppose for every $n\in\naturals$ we have
        $h(n)\in h(\suc{n})$.
    Then $h$ is a natural subsequence selector.
\end{theorem}
\begin{proof}
    Let $C=\{n\in\naturals\mid
        \text{for every $m\in n$ we have $h(m)\in h(n)$}\}$.
    $\emptyset\in C$ by \auto{omega_empty,emptyset}.
    Show that for every $n\in C$ we have $\suc{n}\in C$.
    \begin{subproof}
        Fix $n\in C$.
        Hence $n\in\naturals$ and for every $m\in n$ we have
            $h(m)\in h(n)$.
        $\suc{n}\in\naturals$ by \obv{omega_successor}.
        Show that for every $m\in\suc{n}$ we have
            $h(m)\in h(\suc{n})$.
        \begin{subproof}
            Fix $m\in\suc{n}$.
            $\naturals$ is \in-transitive by
                \obv{naturals_is_transitive_set}.
            Hence $m\in\naturals$ by
                \obv{transitiveset_iff_subseteq,subseteq}.
            Hence $m=n$ or $m\in n$ by \obv{suc_iff}.
            \begin{byCase}
                \caseOf{$m=n$.}
                    Follows by \auto.
                \caseOf{$m\neq n$.}
                    Hence $m\in n$.
                    Hence $h(m)\in h(n)$.
                    $h(m)\in\naturals$ by \obv{funs_type_apply}.
                    $h(n)\in\naturals$ by \obv{funs_type_apply}.
                    $h(\suc{n})\in\naturals$ by
                        \obv{funs_type_apply}.
                    Hence $h(\suc{n})$ is an ordinal by
                        \obv{natural_number_is_ordinal}.
                    $h(n)\in h(\suc{n})$.
                    Follows by \obv{ordinal_transitivity}.
            \end{byCase}
        \end{subproof}
    \end{subproof}
    Hence $C$ is an inductive set by \obv{inductive_set}.
    Hence every natural number is an element of $C$ by
        \obv{omega_minimal}.
    Hence for all $m,n\in\naturals$ we have if $m\in n$, then
        $h(m)\in h(n)$ by \auto.
    Follows by \obv{natural_subsequence_selector}.
\end{proof}

\begin{definition}\label{subsequence}
    $t$ is a subsequence of $s$ in $S$ iff
        $s,t\in\funs{\naturals}{S}$ and there exists $h$ such that
            $h$ is a natural subsequence selector and $t=s\circ h$.
\end{definition}

\begin{proposition}\label{subsequence_types}
    Suppose $t$ is a subsequence of $s$ in $S$.
    Then $s,t\in\funs{\naturals}{S}$.
\end{proposition}
\begin{proof}
    Follows by \obv{subsequence}.
\end{proof}

\begin{proposition}\label{subsequence_selector_exists}
    Suppose $t$ is a subsequence of $s$ in $S$.
    Then there exists $h$ such that
        $h$ is a natural subsequence selector and $t=s\circ h$.
\end{proposition}
\begin{proof}
    Follows by \obv{subsequence}.
\end{proof}

\begin{proposition}\label{subsequence_from_selector}
    Let $s\in\funs{\naturals}{S}$.
    Suppose $h$ is a natural subsequence selector.
    Then $s\circ h$ is a subsequence of $s$ in $S$.
\end{proposition}
\begin{proof}
    $h\in\funs{\naturals}{\naturals}$ by
        \obv{natural_subsequence_selector_type}.
    Hence $s\circ h\in\funs{\naturals}{S}$ by \obv{funs_circ}.
    Follows by \obv{subsequence}.
\end{proof}

\begin{proposition}\label{subsequence_from_selector_apply}
    Let $s\in\funs{\naturals}{S}$.
    Suppose $h$ is a natural subsequence selector.
    Let $n$ be a natural number.
    Then $(s\circ h)(n)=s(h(n))$.
\end{proposition}
\begin{proof}
    $h\in\funs{\naturals}{\naturals}$ by
        \obv{natural_subsequence_selector_type}.
    Hence $s$ and $h$ are functions by \obv{funs_is_function}.
    $\dom{s}=\naturals$ and $\dom{h}=\naturals$ by
        \obv{funs}.
    $\ran{h}\subseteq\naturals$ by \obv{funs_ran}.
    Hence $\ran{h}\subseteq\dom{s}$.
    $n\in\dom{h}$.
    Follows by \obv{circ_apply}.
\end{proof}
